\chapter{Sequential Trading Economies}

\section{Incomplete Financial Market with $L = 1$}
\subsection{CPE when $L = 1$}.
\paragraph{(Diamond)} 
A financial market equilibirum $(x, \pi) \in \mathbb{R}^{(S+1)I}_{+} \times \mathbb{R}_{++}^S$
of an economy with $\text{rank}(A) =  J < S$
is constrained Pareto Efficient when $L= 1$.

\paragraph{Proof.} We prove by contradiction, as the proof of the First Welfare Theorem. 
When $L=1$, the spanning condition becomes 
\begin{align*}
    \begin{bmatrix}
        \vdots \\ 
        y^i_s - w^i_s \\
        \vdots
    \end{bmatrix} \in \text{span}(A)
\end{align*}

Let $x$ be a financial market equilibrium. 
Suppose there exists a $y$ that is preferred by all $i$ (with strict preference for at least one $j$.)
Then, we have 
\begin{align*}
    y^i_0 - w^i_0 + \sum_s \pi_s (y^i_s - w^i_s) & \geq x^i_0 - w^i_0 + \sum_s \pi_s (x^i_s - w^i_s) \\ 
    y^j_0 - w^j_0 + \sum_s \pi_s (y^j_s - w^j_s) & > x^j_0 - w^j_0 + \sum_s \pi_s (x^j_s - w^j_s)
\end{align*}
Summing across all $i$ and invoking the Walras' Law, we have
\begin{align*}
    \sum_i y^i_0 - w^i_0 + \sum_s \pi_s  \sum_i (y^i_s - w^i_s) > 0 
\end{align*}
Since $\pi_s  > 0$, we have $\sum_i (y^i_s - w^i_s) > 0 $, contradicting the feasiblity condition. 
\\
\textit{Q.E.D.}


\section{The Representative Agent Theorem}

\subsection*{Statement}

\begin{tcolorbox}
\textbf{Theorem (Representative Agent).} Suppose all $I$ agents have identical VNM preferences with the same homothetic Bernoulli utility:
\[
U^i(x^i) = u(x_0^i) + \sum_{s=1}^S \mathrm{prob}_s\, u(x_s^i), \quad u(\cdot) \text{ homothetic}, \ \forall i.
\]
If markets are complete ($\mathrm{rank}(A) = S$), then there exists a representative agent with
\[
U^R(x) = U(x), \qquad \omega^R = \sum_{i=1}^I \omega^i,
\]
such that the FME of the single-agent economy has the same prices $(p,q)$ as the original $I$-agent FME. The representative agent's equilibrium allocation is autarchic: $x^R = \omega^R$, $z^R = 0$.
\end{tcolorbox}

\subsection*{Three ingredients (memorize these)}

\begin{enumerate}
    \item \textbf{Identical preferences} $\Rightarrow$ same MRS structure across agents.
    \item \textbf{Homotheticity} $\Rightarrow$ individual demand is linear in wealth: $x^i(p, y^i) = y^i \, x(p, 1)$, so demand aggregates linearly.
    \item \textbf{Complete markets} ($\mathrm{rank}(A) = S$) $\Rightarrow$ equilibrium is Pareto efficient and equalizes MRSs across agents, so only aggregate wealth matters.
\end{enumerate}

\subsection*{Proof sketch}

\textbf{Step 1: Linear aggregation from homotheticity.} By homotheticity, $x^i(p, y^i) = y^i x(p,1)$. Identical preferences give $x^i(p,1) = x(p,1)$. Summing across agents:
\[
\sum_i x^i(p, y^i) = \left(\sum_i y^i\right) x(p,1) = x\!\left(p, \sum_i y^i\right).
\]
The last equality reapplies homotheticity in reverse — aggregate demand equals the demand of a single agent endowed with $\sum_i y^i$.

\medskip
\textbf{Step 2: Walras + spot market clearing.} The budget constraint gives $p_s x_s^i = y_s^i$. Summing and using market clearing $\sum_i x_s^i = \sum_i \omega_s^i$:
\[
\sum_i y_s^i = p_s \sum_i \omega_s^i.
\]

\medskip
\textbf{Step 3: Combine.} Substituting into the aggregated demand:
\[
x\!\left(p_s, p_s \sum_i \omega_s^i\right) = \sum_i \omega_s^i.
\]
This is the equilibrium condition for a \emph{single} agent who owns $\sum_i \omega^i$ and faces prices $p$. Hence the same $p$ supports both equilibria.

\medskip
\textbf{Step 4: Where completeness enters.} Complete markets ensure the original $I$-agent equilibrium is Pareto efficient, so all agents share a common pricing kernel $m$. This is what makes the aggregation valid \emph{across states} (not just within a state) and pins down $q$ from the representative agent's MRS: $q = \mathbb{E}[m^R A]$.

\subsection*{Why it matters}

The theorem reduces equilibrium price computation to a \emph{single-agent} problem: just solve
\[
q_j = \sum_{s} \mathrm{prob}_s \frac{u'\!\left(\sum_i \omega_s^i\right)}{u'\!\left(\sum_i \omega_0^i\right)} a_{sj},
\]
plugging in aggregate endowments. No need to solve for the individual allocation $(x^i, z^i)_{i=1}^I$. This is the foundation of essentially all macro/finance asset-pricing models (CCAPM, Lucas tree, etc.).

\subsection*{Memory hooks}

\begin{itemize}
    \item \textbf{Slogan:} ``Identical + homothetic + complete $\Rightarrow$ aggregate.''
    \item \textbf{Each assumption kills one obstruction:} identical preferences kill heterogeneous demand functions; homotheticity kills wealth-distribution effects; completeness kills MRS dispersion across agents.
    \item \textbf{Compare to uniqueness Proposition 6(ii):} the static version (identical + homothetic) gives uniqueness in a pure exchange economy. The representative agent theorem is the dynamic/state-contingent analogue, with completeness added to handle the extra dimension.
    \item \textbf{Equilibrium is autarchic:} $x^R = \omega^R$, $z^R = 0$ — the representative agent trades with herself, so net trades vanish.
\end{itemize}
